\chapter{Cartan fields, boundary currents and gravitational states}
\label{mht:ch:gravitational-models}

A boundary stress tensor supplies a transformation law and a central
charge. A gravitational comparison must also supply bulk fields,
their gauge symmetry, an action and a state pairing. In three
dimensions the classical field comparison is unusually explicit:
an invertible coframe converts a flat Cartan connection into an
Einstein metric. This conversion determines neither a quantum
integration cycle nor a boundary state.

We begin with the action and its infinitesimal complexes. The
boundary matrix calculation then fixes the classical central term.
Finite Gaussian integrals, Virasoro states and marked parallel
transport give further comparisons with specified domains. Together
they identify the data needed to extend the classical correspondence
to the quantum comparisons of Chapter~\ref{mht:ch:compact-gravity}.

\section{The coframe and the action}

Let $M$ be an oriented three-manifold, and let $V\to M$ be an
oriented rank-three bundle with metric $\eta$ of signature $(-,+,+)$.
A coframe is a bundle isomorphism $e:TM\to V$.
It defines $g=e^*\eta$.
Choose $\ell>0$ and set $\Lambda=-\ell^{-2}$.
Let $G>0$ denote Newton's constant.

In an orthonormal frame, put $\epsilon_{012}=1$ and raise indices
with $\eta$. The Lorentz algebra has generators $J_a$ satisfying
$[J_a,J_b]=\epsilon_{ab}{}^cJ_c$.

Write the spin connection as $\omega=\omega^aJ_a$ and its curvature as
\[
 F_\omega=d\omega+\tfrac12[\omega,\omega]=R^aJ_a.
\]
Brackets of Lie-valued forms include their exterior product.
The invariant bilinear form used in each chiral sector is
\[
 B(J_a,J_b)=\tfrac12\eta_{ab}.
\]
It is the fundamental matrix trace in the usual real
$\mathfrak{sl}_2$ realization. On a chosen trivialization define
\[
 \operatorname{CS}_B(A)
 =B(A\wedge dA)+\tfrac13B(A\wedge[A,A]),
 \qquad A^\pm=\omega\pm\ell^{-1}e.
\]
Here $e$ is identified with a Lorentz-algebra-valued one-form by
the oriented metric on $V$.

\begin{proposition}\label{mht:prop:gravity-action-normalization}
The two Chern--Simons forms satisfy
\begin{equation}\label{mht:eq:gravity-cs-difference}
 \begin{split}
 \operatorname{CS}_B(A^+)-\operatorname{CS}_B(A^-)
 ={}&\frac4\ell B(e\wedge F_\omega)
       +\frac{2}{3\ell^3}B(e\wedge[e,e])\\
    &-\frac2\ell dB(\omega\wedge e).
 \end{split}
\end{equation}
Consequently, for $k=\ell/(4G)$, the bulk part of
$\frac{k}{4\pi}\int_M(\operatorname{CS}_B(A^+)-
\operatorname{CS}_B(A^-))$ is
\begin{equation}\label{mht:eq:gravity-palatini}
 \frac1{8\pi G}\int_M e_a\wedge R^a
 +\frac1{48\pi G\ell^2}\int_M
       \epsilon_{abc}e^a\wedge e^b\wedge e^c.
\end{equation}
Its displayed boundary term is
$-\frac{k}{4\pi\ell}\int_{\partial M}\eta_{ab}\omega^a\wedge e^b$.
\end{proposition}

\begin{proof}
Set $a=e/\ell$ and expand in $a$.
The quadratic terms contribute
$2B(\omega\wedge da)+2B(a\wedge d\omega)$.
Since
$dB(\omega\wedge a)=B(d\omega\wedge a)-B(\omega\wedge da)$,
their sum is $4B(a\wedge d\omega)-2dB(\omega\wedge a)$.

Invariance of $B$ makes the three terms containing one $a$ in
the cubic expression equal. Their sum is
$2B(a\wedge[\omega,\omega])$.
The remaining odd term is $\frac23B(a\wedge[a,a])$.
These terms give \eqref{mht:eq:gravity-cs-difference}.

The identities
$B(e\wedge F_\omega)=\frac12e_a\wedge R^a$ and
$B(e\wedge[e,e])=\frac12\epsilon_{abc}e^a\wedge e^b\wedge e^c$
give \eqref{mht:eq:gravity-palatini} after substituting $k$.
\end{proof}

With $R^a=\frac12\epsilon^a{}_{bc}R^{bc}$, the first term
in \eqref{mht:eq:gravity-palatini} is
$\frac1{16\pi G}\int\epsilon_{abc}e^a\wedge R^{bc}$.
This is the first-order Einstein--Hilbert normalization.
The trace convention and the relation $k=\ell/(4G)$ are the
ones used in the Chern--Simons description of three-dimensional
gravity.\footnote{See Campoleoni, Fredenhagen, Pfenninger and Theisen,
\emph{Asymptotic symmetries of three-dimensional gravity coupled to
higher-spin fields}, Section~2.2, equations~(2.22)--(2.25),
\url{https://arxiv.org/abs/1008.4744}.}
A boundary variational problem must further specify its fixed
fields and any additional boundary action.

The Cartan algebra $\mathfrak g_\Lambda$ has generators $J_a,P_a$ with
\[
 [J_a,P_b]=\epsilon_{ab}{}^cP_c,
 \qquad [P_a,P_b]=\ell^{-2}\epsilon_{ab}{}^cJ_c,
 \qquad \mathcal A=\omega^aJ_a+e^aP_a.
\]
For the off-diagonal pairing $b(J_a,P_b)=\eta_{ab}$, put
$J_a^\pm=(J_a\pm\ell P_a)/2$.
Then $b(J_a^\pm,J_b^\pm)=\pm\ell B(J_a,J_b)$ and
$b(J_a^+,J_b^-)=0$.
Thus the Cartan action is
$\frac{k}{4\pi\ell}\int_M\operatorname{CS}_b(\mathcal A)$.
Its level is $k/\ell=1/(4G)$ when the two sector levels
are $k$ and $-k$.

\begin{proposition}\label{mht:prop:gravity-flat-coframe}
On the invertible-coframe locus, $F_{\mathcal A}=0$ is equivalent
to the torsion-free Einstein equation with cosmological constant
$\Lambda=-\ell^{-2}$.
We use $G_{\mu\nu}=\operatorname{Ric}_{\mu\nu}
-\frac12Rg_{\mu\nu}$ for the Einstein tensor.
\end{proposition}

\begin{proof}
The translation component of the curvature is
$T^a=d e^a+\epsilon^a{}_{bc}\omega^b\wedge e^c$.
Its vanishing determines the metric-compatible torsion-free
connection because $e$ is invertible.
The rotation component is
$R^a+(2\ell^2)^{-1}\epsilon^a{}_{bc}e^b\wedge e^c$.
Its vanishing says that the metric has constant curvature
$-\ell^{-2}$.
In dimension three the Riemann tensor is determined by its Ricci
contraction, so this is equivalent to
$G_{\mu\nu}+\Lambda g_{\mu\nu}=0$.
The same identities prove the converse.
\end{proof}

Invertibility is essential: $e=0$ and a flat $\omega$ solve the
Cartan equations but give a degenerate metric.
The classical comparison therefore applies to an open field domain.
It does not identify a path integral over all connections with one
over nondegenerate metrics.\footnote{Witten,
\emph{$2+1$ dimensional gravity as an exactly soluble system},
Section~2, explicitly separates the invertible-coframe argument from
the larger gauge-field domain. See
\url{https://doi.org/10.1016/0550-3213(88)90143-5}.}

\section{Eliminating the algebraic pairs}

Fix a flat background $\mathcal A$ with invertible coframe.
Its unshifted deformation complex has degrees zero through three:
\begin{equation}\label{mht:eq:gravity-flat-complex}
 \Omega^0(M,\mathfrak g_\Lambda)
 \xrightarrow{d_{\mathcal A}}
 \Omega^1(M,\mathfrak g_\Lambda)
 \xrightarrow{d_{\mathcal A}}
 \Omega^2(M,\mathfrak g_\Lambda)
 \xrightarrow{d_{\mathcal A}}
 \Omega^3(M,\mathfrak g_\Lambda).
\end{equation}
Flatness gives $d_{\mathcal A}^2=[F_{\mathcal A},-]=0$.

Write a parameter as $(\xi,\lambda)$, a fluctuation as
$(u,v)=(\delta e,\delta\omega)$, and a curvature variation as
$(t,r)$ in the translation and rotation summands.
Put $D_\omega=d+[\omega,-]$.
The differentials are
\begin{align}
 (\xi,\lambda)&\longmapsto
 (D_\omega\xi+[e,\lambda],
  D_\omega\lambda+\ell^{-2}[e,\xi]),\label{mht:eq:gravity-gauge}\\
 (u,v)&\longmapsto
 (D_\omega u+[e,v],
  D_\omega v+\ell^{-2}[e,u]),\label{mht:eq:gravity-curvature}\\
 (t,r)&\longmapsto
 (D_\omega t+[e,r],
  D_\omega r+\ell^{-2}[e,t]).\label{mht:eq:gravity-bianchi}
\end{align}
They are the components of the same flat differential.

We shall remove three pairs connected by invertible algebraic
operators. The relevant cancellation also specifies the comparison
maps.

\begin{lemma}\label{mht:lem:gravity-cancellation}
Suppose a differential between two consecutive terms of a complex
has block form
$\left(\begin{smallmatrix}a&b\\c&d\end{smallmatrix}\right)$
on $U\oplus A\to V\oplus B$, with $a:U\to V$ invertible.
The pair $U\to V$ contracts, leaving differential
$d-ca^{-1}b:A\to B$.
The lift of $x\in A$ in that degree is $(-a^{-1}bx,x)$,
and projection in the next degree sends $(v,y)$ to $y-ca^{-1}v$.
\end{lemma}

\begin{proof}
Multiply the differential on the right and left by
\[
 \begin{pmatrix}1&-a^{-1}b\\0&1\end{pmatrix},
 \qquad
 \begin{pmatrix}1&0\\-ca^{-1}&1\end{pmatrix},
\]
respectively. It becomes $a\oplus(d-ca^{-1}b)$.
The square-zero identity forces the adjacent differentials to have
zero components into and out of this isolated invertible block.
Its inverse is a contracting homotopy.
Undoing the two basis changes gives the stated lift and projection.
\end{proof}

Use $e$ to regard $u$ as an endomorphism of $TM$.
Its symmetric part is the metric variation
\[
 h_{\mu\nu}=\eta_{ab}
       (e^a_\mu u^b_\nu+e^a_\nu u^b_\mu).
\]
The Lorentz parameter maps isomorphically to the skew part of $u$.
Cancelling this degree-zero/degree-one pair removes the frame
rotation and retains $h$.

Next consider the torsion map on a connection variation,
$L_e(v)=v^a{}_b\wedge e^b$.
In an orthonormal coframe, write
$v_{ab}=v_{abc}e^c$, where $v_{abc}=-v_{bac}$, and
$L_e(v)^a=\frac12t^a{}_{bc}e^b\wedge e^c$.
Then
\begin{equation}\label{mht:eq:gravity-torsion-inverse}
 t_{abc}=v_{acb}-v_{abc},
 \qquad
 v_{abc}=\tfrac12(t_{cab}-t_{abc}-t_{bca}).
\end{equation}
This is an isomorphism between two rank-nine bundles.
It cancels $v$ against the torsion component in the next degree.
The lift of $u$ is the unique $v$ satisfying
$L_e(v)=-D_\omega u$.

After this cancellation, the translation Bianchi equation is
$[e,r]=0$.
Dualizing the form indices identifies $r$ with an endomorphism
of $V$. The map $r\mapsto[e,r]$ identifies its skew part with
a rank-three Bianchi component.
Cancelling this final pair leaves a symmetric curvature tensor
and the divergence identity.

\begin{theorem}\label{mht:thm:gravity-metric-complex}
The three cancellations give a homotopy equivalence, by local
differential operators, from \eqref{mht:eq:gravity-flat-complex} to
\begin{equation}\label{mht:eq:gravity-metric-complex}
 \Gamma(TM)\xrightarrow{K}\Gamma(S^2T^*M)
 \xrightarrow{E_\Lambda}\Gamma(S^2T^*M)
 \xrightarrow{\operatorname{div}}\Gamma(T^*M).
\end{equation}
Here $K\zeta=\mathcal L_\zeta g$, $E_\Lambda$ is the variation
of $G+\Lambda g$, and the final map is covariant divergence.
On a sufficiently small geodesically convex open set, this complex
resolves the sheaf of Killing fields.
\end{theorem}

\begin{proof}
Each cancellation above has a pointwise inverse, so
Lemma~\ref{mht:lem:gravity-cancellation} supplies chain maps and
homotopies. The remaining parameter is the vector field
$\zeta=e^{-1}\xi$.
Cartan's identity
$\mathcal L_\zeta\mathcal A=d_{\mathcal A}(\iota_\zeta\mathcal A)$
on a flat background shows that its metric variation is $K\zeta$.
A compensating Lorentz rotation gives the symmetric-frame lift
selected by the first cancellation.

Solving the torsion equation gives the variation of the
Levi--Civita connection. The rotation curvature equation therefore
becomes the variation of $G+\Lambda g$.
Its remaining Bianchi component is covariant divergence.
The explicit formulas below identify these two operators.
Thus the cancelled complex is \eqref{mht:eq:gravity-metric-complex}.

A flat connection has a local parallel trivialization.
The ordinary Poincar\'e homotopy then contracts its positive-degree
de Rham complex on a small convex open set.
The degree-zero kernel transfers to $K\zeta=0$ under the chain
equivalence. This gives the asserted local resolution.
\end{proof}

The operator $E_\Lambda$ can be written without commuting covariant
derivatives. Put $h=g^{\mu\nu}h_{\mu\nu}$ and define
\begin{align}
 C^\rho{}_{\mu\nu}
 &=\tfrac12g^{\rho\sigma}
   (\nabla_\mu h_{\nu\sigma}+\nabla_\nu h_{\mu\sigma}
     -\nabla_\sigma h_{\mu\nu}),\label{mht:eq:gravity-christoffel-variation}\\
 r_{\mu\nu}
 &=\nabla_\rho C^\rho{}_{\mu\nu}
       -\nabla_\nu C^\rho{}_{\mu\rho},\qquad
 r=g^{\mu\nu}r_{\mu\nu}-2\Lambda h.
\end{align}
Since $\operatorname{Ric}=2\Lambda g$ and $R=6\Lambda$,
\begin{equation}\label{mht:eq:gravity-einstein-variation}
 (E_\Lambda h)_{\mu\nu}
   =r_{\mu\nu}-\tfrac12g_{\mu\nu}r-2\Lambda h_{\mu\nu}.
\end{equation}
Diffeomorphism covariance gives $E_\Lambda K=0$.
Varying the contracted Bianchi identity at the background solution
gives $\nabla^\mu(E_\Lambda h)_{\mu\nu}=0$.

A constant rescaling checks the curvature terms:
for $h_{\mu\nu}=g_{\mu\nu}$, one has $C=0$, $r=-6\Lambda$
and $E_\Lambda h=\Lambda g$.
Equivalently, $G((1+t)g)=G(g)$ as a covariant tensor, while
the cosmological term varies by $t\Lambda g$.
The kernel of $(u,v)\mapsto h$ consists of Lorentz variations
only after imposing the linearized torsion equation.
Without that condition, $(0,v)$ is in the kernel for every $v$.

For the Calabi presentation, use the curvature defect
\[
 \mathcal R_\Lambda(g)_{\rho\mu\sigma\nu}
 =R_{\rho\mu\sigma\nu}
  -\Lambda(g_{\rho\sigma}g_{\mu\nu}
             -g_{\rho\nu}g_{\mu\sigma}).
\]
It vanishes at the background, so its variation kills $K\zeta$.

In dimension three, Ricci contraction followed by subtraction of
half the scalar trace is an isomorphism from algebraic curvature
tensors to symmetric two-tensors.
Applied to $\mathcal R_\Lambda(g)$, it gives $G+\Lambda g$.
Its linearization therefore identifies the corrected Calabi
curvature operator with \eqref{mht:eq:gravity-einstein-variation}.
This comparison concerns differential complexes on open sets.
Boundary conditions and Gaussian measures require additional maps.

\section{The projective boundary calculation}

Consider the matrix connection $\partial_z+a$ on a coordinate
interval or formal disc, where
\[
 a=\begin{pmatrix}0&q\\1&0\end{pmatrix}.
\]
An infinitesimal gauge parameter $\lambda$ acts by
$\delta a=\partial_z\lambda+[a,\lambda]$.
Preserving the lower-left entry and the zero diagonal leaves one
function $\varepsilon$:
\begin{equation}\label{mht:eq:gravity-ds-parameter}
 \lambda_\varepsilon=
 \begin{pmatrix}
  -\varepsilon'/2&q\varepsilon-\varepsilon''/2\\
  \varepsilon&\varepsilon'/2
 \end{pmatrix},
 \qquad
 \delta_\varepsilon q
 =\varepsilon q'+2\varepsilon'q-\tfrac12\varepsilon'''.
\end{equation}
Primes denote $z$-derivatives.
Indeed the lower-left equation gives the diagonal entries, and
the diagonal equation determines the upper-right entry.
Substitution gives the displayed variation of $q$.

Choose the chiral orientation for which the charge variation is
\[
 \delta Q_\varepsilon
   =-\frac{k}{2\pi}\int\operatorname{tr}
                       (\lambda_\varepsilon\delta a)\,dz.
\]
Since the trace is $\varepsilon\delta q$, the stress density
in the convention $Q_\varepsilon=(2\pi)^{-1}\int\varepsilon T\,dz$
is $T=-kq$. Consequently
\begin{equation}\label{mht:eq:gravity-classical-stress}
 \delta_\varepsilon T
 =\varepsilon T'+2\varepsilon'T+\frac{k}{2}\varepsilon'''.
\end{equation}
The Virasoro central term $c\varepsilon'''/12$ therefore has
$c=6k=3\ell/(2G)$.
This is the classical Brown--Henneaux normalization for the
specified boundary reduction.\footnote{The boundary constraints,
charges and central term are derived in Campoleoni et al.,
Section~4.2, equations~(4.13)--(4.26),
\url{https://arxiv.org/abs/1008.4744}.}

The same law follows from the scalar equation $\psi''=q\psi$.
For $z=f(w)$ and a local square root of $f'$, set
$\widetilde\psi(w)=(f'(w))^{-1/2}\psi(f(w))$.
Two differentiations give
\begin{equation}\label{mht:eq:gravity-projective-change}
 \widetilde q(w)=(f')^2q(f(w))-\tfrac12\{f,w\},
 \qquad
 \{f,w\}=\frac{f'''}{f'}-\frac32\left(\frac{f''}{f'}\right)^2.
\end{equation}

The Schwarzian chain rule makes this a coordinate-gluing law.
It defines a projective connection, whose differences are
quadratic differentials.
Taking $f(w)=w+t\varepsilon(w)$ recovers
\eqref{mht:eq:gravity-ds-parameter}.
The solution half-density uses a chosen square root locally.
The projective coefficient itself is independent of its sign.

For a holomorphic projective background, the infinitesimal action
gives the two-term Dolbeault complex
\[
 \Omega^{0,\bullet}(X,T_X)
 \xrightarrow{\ \varepsilon\mapsto
       \varepsilon T'+2\varepsilon'T+(c/12)\varepsilon'''\ }
 \Omega^{0,\bullet}(X,K_X^2).
\]
A family of mixed boundary operations over these backgrounds also
needs coordinate-equivariance for every operation and coherence map.
Choose a differential graded Lie model of this action and a
compatible action on the relative deformation algebra.
Their semidirect product describes the coupled deformation problem.
A conformal vector specifies the stress action on its vertex algebra.
Compatibility with additional mixed operations is further data.

\subsection{The quantum affine parameter}

Let $\kappa\in\mathbb C\setminus\{-2\}$ be the level of the ordinary affine
$\mathfrak{sl}_2$ vertex algebra with trace pairing $B$.
Its Sugawara central charge is $3\kappa/(\kappa+2)$.
Principal quantum Drinfeld--Sokolov reduction adds the derivative
of the Cartan current to the stress tensor.
It also adds one fermionic ghost pair of weights zero and one.

The Cartan current has double-pole coefficient $\kappa/2$.
Its derivative changes the central charge by $-6\kappa$, and
the ghost pair contributes $-2$.
Thus the algebraic reduced central charge is
\begin{equation}\label{mht:eq:gravity-quantum-ds}
 c_{\mathrm{DS}}(\kappa)
 =\frac{3\kappa}{\kappa+2}-6\kappa-2
 =1-\frac{6(\kappa+1)^2}{\kappa+2}.
\end{equation}

The free-field formula makes the normalization explicit.
For $b(z)b(w)\sim(z-w)^{-2}$ put
\[
 T_\nu=\tfrac12:b^2:+\nu\partial b,
 \qquad \nu=\frac{\kappa+1}{\sqrt{2(\kappa+2)}}.
\]
Choose either square root in this expression.
Wick contraction gives the identity coefficient
$\frac12-6\nu^2$ in the fourth-order pole of $T_\nu(z)T_\nu(w)$.
The other contractions give $2T_\nu$ and $\partial T_\nu$ in
the second- and first-order poles.
Hence its central charge is $1-12\nu^2$, equal to
\eqref{mht:eq:gravity-quantum-ds}.

This is the standard free-field realization of the reduced
$\mathfrak{sl}_2$ stress tensor.\footnote{Arakawa,
\emph{Introduction to $W$-algebras and their representation theory},
Section~5.4 and Example~5.22,
\url{https://arxiv.org/abs/1605.00138}.}

The affine parameter $\kappa$ and the gravitational coefficient
$k$ have distinct definitions. Their quantum relation must be
derived from the physical current and ghost insertion, its finite
terms and its integration convention.
Equation~\eqref{mht:eq:gravity-classical-stress} fixes a classical
charge normalization. Equation~\eqref{mht:eq:gravity-quantum-ds}
is an identity in the specified algebraic reduction.

\section{Principal embeddings and holonomy}

Let $V_{N-1}=\operatorname{Sym}^{N-1}\mathbb C^2$ and let
$\mathfrak{sl}_2$ act on $\operatorname{End}(V_{N-1})$ by
commutators. The defining action gives the principal embedding
in $\mathfrak{sl}_N$.
If $E,F,H$ are its standard generators, $E^j$ is a highest-weight
vector of weight $2j$ for $0\le j\le N-1$.
Its iterated commutators with $F$ span a module of dimension $2j+1$.
The raising and lowering relation is
\[
 (\operatorname{ad}E)(\operatorname{ad}F)^rE^j
 =r(2j-r+1)(\operatorname{ad}F)^{r-1}E^j.
\]
The finite-dimensional highest-weight relations give $V_{2j}$.

These modules have distinct highest weights, and their dimensions
sum to $N^2$.
Removing the scalar module gives
\begin{equation}\label{mht:eq:gravity-principal-decomposition}
 \mathfrak{sl}_N=\bigoplus_{j=1}^{N-1}V_{2j},
 \qquad \sum_{j=1}^{N-1}(2j+1)=N^2-1.
\end{equation}

In a supplied $\mathfrak{sl}_N\oplus\mathfrak{sl}_N$
Chern--Simons theory, each summand gives one pair of one-form
fields in the Lorentz representation $V_{2j}$.
These are the frame variables of spin $s=j+1$.
Their linearized torsion and curvature equations on an AdS background
are the spin-$s$ versions of \eqref{mht:eq:gravity-curvature}.
Eliminating the spin connection and the Lorentz-type gauge component
gives the corresponding Fronsdal field with its trace constraints.
This field interpretation uses the invertible background coframe
and the supplied action.\footnote{The torsion solution, metric
field and gauge transformation are given explicitly in
Campoleoni et al., Section~2.2, equations~(2.28)--(2.37),
\url{https://arxiv.org/abs/1008.4744}.}

Principal Drinfeld--Sokolov boundary reduction has currents of
weights $2,\ldots,N$, forming the classical $W_N$ algebra.
The representation match identifies the possible spin labels.
Realizing a specified pair of $W_N$ boundary theories still requires
their boundary action, real form, states and map from these bulk
fields.

The distinction between a connection and its exponentiated holonomy
is already visible for constant matrices.
For example, the companion connection
\[
 a=\begin{pmatrix}0&0&w_3\\1&0&w_2\\0&1&0\end{pmatrix}
\]
has characteristic polynomial $\lambda^3-w_2\lambda-w_3$ and
\[
 \operatorname{tr}(a^2)=2w_2,
 \qquad \operatorname{tr}(a^3)=3w_3.
\]
These are polynomial coordinates on the constant companion slice.

Parallel transport for $\partial+a$ around a circle of length $L$
is $\exp(-La)$.
For $a=\operatorname{diag}(x,-x)$, replacing $x$ by
$x+2\pi i n/L$ leaves this holonomy unchanged and generally
changes $\operatorname{tr}(a^2)$.
Thus recovery of the constant charges from holonomy needs a
logarithm branch and the relevant gauge identifications.
For nonconstant connections it also needs the ordered transport.

A smooth flat connection on a contractible disk has identity
holonomy around its boundary in its actual gauge group.
A central element is indistinguishable from the identity only
after passage to a group quotient that kills it.
Thermal regularity, the choice of quotient and the specified
state sector therefore enter a higher-spin filling problem.

\section{Stress tensors from an action}

A stress tensor is obtained by varying a supplied matter action.
Let $S_{\mathrm m}(g,\varphi)$ be invariant under compactly supported
diffeomorphisms in the interior of $M$.
Define $T^{\mu\nu}$ and the matter Euler--Lagrange expression
$\mathcal E_\varphi$ by
\[
 \delta S_{\mathrm m}
 =\int_M\mathrm{vol}_g
   \left(\tfrac12T^{\mu\nu}\delta g_{\mu\nu}
                  +\mathcal E_\varphi\delta\varphi\right).
\]
This equation fixes the sign convention for $T$.

For a scalar field, diffeomorphism invariance implies
\begin{equation}\label{mht:eq:gravity-noether-stress}
 \nabla_\mu T^{\mu\nu}=\mathcal E_\varphi\nabla^\nu\varphi.
\end{equation}
Indeed, substitute
$\delta g_{\mu\nu}=\nabla_\mu\zeta_\nu+\nabla_\nu\zeta_\mu$
and $\delta\varphi=\zeta^\mu\partial_\mu\varphi$, then integrate
the metric term by parts. Since $\zeta$ is arbitrary,
\eqref{mht:eq:gravity-noether-stress} follows.
Stress conservation holds on the matter equations of motion.
With a boundary, the same variation also determines a boundary flux.

The full metric dependence matters for a coupled action.
For the Euclidean scalar kinetic density, a symmetric variation $h$ of the
identity metric gives
\begin{equation}\label{mht:eq:gravity-quadratic-density}
 \begin{split}
 \sqrt{\det(I+th)}(I+th)^{-1}
 ={}&I+t\left(\tfrac12\operatorname{tr}(h)I-h\right)\\
 &+t^2\left[
 \left(\frac{\operatorname{tr}(h)^2}{8}
        -\frac{\operatorname{tr}(h^2)}4\right)I
 -\frac{\operatorname{tr}(h)}2h+h^2\right]+O(t^3).
 \end{split}
\end{equation}
This follows by multiplying the determinant and inverse-matrix
series. For $h=I$ in dimension three, the result is
$(1+t)^{1/2}I$, with nonzero quadratic coefficient $-I/8$.
The linear stress coupling is therefore only the first variation
of this action.

There is a precise finite Hamiltonian statement behind the
equivariance condition. Let $(P,\{-,-\})$ be a symplectic manifold.
Let $\mathfrak h$ be a finite-dimensional Lie algebra with basis
$(t_a)$ and bracket $[t_a,t_b]=f_{ab}{}^ct_c$.
Choose functions $\mu_a$ on $P$.

Adjoin odd ghosts $c^a$ of degree one and odd conjugates $b_a$
of degree minus one, with even Poisson bracket
$\{b_a,c^b\}=\delta_a^b$.
The functions on $P$ Poisson-commute with these new coordinates.
Define the degree-one charge
\[
 \Omega=c^a\mu_a-\tfrac12f_{ab}{}^c c^ac^bb_c.
\]

\begin{proposition}\label{mht:prop:gravity-hamiltonian-charge}
The charge satisfies $\{\Omega,\Omega\}=0$ precisely when the moment functions
are equivariant:
\[
 \tfrac12\{\Omega,\Omega\}
 =\tfrac12c^ac^b
       \bigl(\{\mu_a,\mu_b\}-f_{ab}{}^c\mu_c\bigr).
\]
\end{proposition}

\begin{proof}
The bracket of the two terms linear in $c$ gives
$c^ac^b\{\mu_a,\mu_b\}$.
The two cross terms give $-f_{ab}{}^c c^ac^b\mu_c$.
The term with three ghosts and one conjugate is the Jacobi
expression $f_{[ab}{}^d f_{c]d}{}^e$, which vanishes.
The independent exterior monomials $c^ac^b$ give the equivalence.
\end{proof}

For example, on $\mathbb R^2$ with $\{q,p\}=1$, the functions
$\mu_E=q^2/2$, $\mu_F=-p^2/2$ and $\mu_H=-qp$
satisfy the $\mathfrak{sl}_2$ relations.
They give a concrete square-zero charge.
The construction starts with a symplectic structure and moment
functions. Coordinate covariance by itself supplies neither.

A supplied complete $L_\infty$ action similarly defines a
semidirect deformation algebra. If its local realization and
dual pairing are given, its shifted cotangent action has the
form $S(\alpha,\beta)=\langle\beta,Q(\alpha)\rangle$.
The identity $Q^2=0$ implies the classical master equation
for this cotangent action.
Its independent dual fields are part of its definition.
Comparison with gravitational BV fields must preserve the action,
the BV pairing and the boundary terms on the stated coframe domain.

\section{Gaussian determinants and their measures}

The preceding chain equivalence preserves deformation cohomology.
A Gaussian integral additionally uses a quadratic form and a
measure. We first compute it in finite dimensions, where these
data can all be fixed.

Let
$0\to C^0\xrightarrow{d_0}C^1\xrightarrow{d_1}C^2
\xrightarrow{d_2}C^3\to0$
be an acyclic finite-dimensional real Hilbert complex.
Set
\[
 \Delta_j=d_{j-1}d_{j-1}^*+d_j^*d_j,
 \qquad
 \tau(C)=\prod_{j=0}^3(\det\Delta_j)^{(-1)^{j+1}j/2}.
\]
Absent differentials are zero.
Suppose Hodge duality identifies the spectra of $\Delta_3$
with $\Delta_0$ and of $\Delta_2$ with $\Delta_1$.
On $H=(\operatorname{im}d_0)^\perp\subset C^1$, suppose the
quadratic operator $D$ is invertible and self-adjoint, with
$D^2=\Delta_1|_H$.
These are the finite data corresponding to a gauge-fixed
three-dimensional Chern--Simons quadratic form.

Use the Euclidean measures of the given inner products for
fields and gauge parameters. Their quotient contributes
the orbit Jacobian $(\det\Delta_0)^{1/2}$.
For $k\in\mathbb R^\times$, define the oscillatory integral by a small positive
Gaussian damping, with bosonic measure $(2\pi)^{-\dim H/2}dx$.

\begin{proposition}\label{mht:prop:gravity-finite-torsion}
With $r=\dim H$, this normalized quotient integral is
\begin{equation}\label{mht:eq:gravity-finite-gaussian}
 Z_k(C,D)=|k|^{-r/2}\tau(C)^{1/2}
       \exp\left(\frac{\pi i}{4}\operatorname{sign}(kD)\right).
\end{equation}
Here $\operatorname{sign}(kD)$ is the number of positive
eigenvalues minus the number of negative eigenvalues.
\end{proposition}

\begin{proof}
Orthogonal decomposition separates the gauge orbit
$\operatorname{im}d_0$ from $H$.
The linear map from gauge parameters to the orbit has Jacobian
$(\det d_0^*d_0)^{1/2}=(\det\Delta_0)^{1/2}$.
Diagonalizing $D$ reduces the remaining integral to the
one-dimensional damped Fresnel formula.
It gives
\[
 |k|^{-r/2}(\det\Delta_0)^{1/2}
        (\det\Delta_1|_H)^{-1/4}
 \exp\left(\frac{\pi i}{4}\operatorname{sign}(kD)\right).
\]

The exact part of $C^1$ has the nonzero spectrum of $\Delta_0$,
so $\det\Delta_1=(\det\Delta_0)(\det\Delta_1|_H)$.
Hodge duality in the definition of $\tau$ gives
$\tau=(\det\Delta_0)^{3/2}(\det\Delta_1)^{-1/2}$.
These identities give \eqref{mht:eq:gravity-finite-gaussian}.
\end{proof}

Even acyclic complexes can have different torsions.
For $a,b>0$, take dimensions $(1,2,2,1)$ and differentials
\[
 d_0=\begin{pmatrix}a\\0\end{pmatrix},\qquad
 d_1=\begin{pmatrix}0&0\\0&b\end{pmatrix},\qquad
 d_2=\begin{pmatrix}a&0\end{pmatrix}.
\]
With the standard inner products, $\tau=a^2/b$.
Changing $a$ changes the Gaussian magnitude, although every
complex remains acyclic and the determinant-line orientations
can be kept fixed.
An abstract quasi-isomorphism therefore does not determine the
metric torsion or the Gaussian normalization.

For a closed Riemannian three-manifold and an acyclic unitary flat
bundle, the analytic counterpart uses zeta determinants of the
elliptic de Rham Laplacians.
The modulus of the gauge-fixed one-loop expression has the same
determinant exponents. Its phase uses an eta invariant and a
chosen framing correction.\footnote{Witten,
\emph{Quantum field theory and the Jones polynomial},
Section~2, pages~356--359, derives the gauge-fixed determinants
and their phase. The formula here fixes the torsion convention
explicitly in \eqref{mht:eq:gravity-finite-gaussian}.}

These analytic hypotheses differ from those of a noncompact
gravitational real form on a solid torus with boundary.
That problem also requires closed operator domains, a specified
integration cycle, residual-mode measures and normalized cut states.

The chain maps in Theorem~\ref{mht:thm:gravity-metric-complex}
must transport these Gaussian data.
The Jacobians of their cancelled pairs also enter the determinant
comparison.
The source-dependent finite terms and the determinant-line
connection must then agree under the permitted local corrections.
Equality of the local deformation complexes is the classical
part of this comparison.

Algebraic cancellation of curvature has its own precise meaning.
Let $(\mathfrak g,d,[-,-],\kappa)$ be a curved differential graded
Lie algebra with central closed curvature $\kappa$ of degree two.
Adjoin a central generator $u$ of degree one with $du=\kappa$.
The extended differential squares to zero, and
\[
 \kappa+d(-u)+\tfrac12[-u,-u]=0.
\]
Twisting by $-u$ therefore gives zero curvature.
This algebraic construction does not supply an invertible field theory.

Geometrically, opposite line curvatures make the tensor connection
flat. Trivializing it also requires its holonomy and sewing data.
For example, the connection $d+ia\,d\theta$ on the trivial line
over a circle has zero curvature and holonomy $e^{-2\pi ia}$.
At $a=1/2$ its holonomy is $-1$.
Cancellation by an invertible theory thus requires the complete
inverse theory and its compatible boundary maps.

\section{Virasoro states and analytic continuation}

State comparisons can be tested before a full Fock construction.
Let $c\in\mathbb R$ and let the Virasoro generators satisfy
\[
 [L_m,L_n]=(m-n)L_{m+n}
       +\frac{c}{12}(m^3-m)\delta_{m+n,0},
\]
and normalize a vacuum $v$ by
$L_nv=0$ for $n\ge-1$ and $\langle v,v\rangle=1$.
The contravariant Hermitian form obeys $L_n^*=L_{-n}$.
No positivity at all weights is assumed.

\begin{proposition}\label{mht:prop:gravity-vacuum-gram}
The forms at weights two and three are $(c/2)$ and $(2c)$.
At weight four, in the basis $L_{-4}v,L_{-2}^2v$, the form is
\begin{equation}\label{mht:eq:gravity-gram-four}
 G_4=\begin{pmatrix}5c&3c\\3c&c^2/2+4c\end{pmatrix},
 \qquad \det G_4=\frac{c^2(5c+22)}2.
\end{equation}
For $c>0$, these displayed forms are positive.
\end{proposition}

\begin{proof}
The first two norms are the vacuum values of
$[L_2,L_{-2}]$ and $[L_3,L_{-3}]$.
At weight four,
\[
 L_4L_{-2}^2v=3cv,
 \qquad L_2L_{-2}^2v=(8+c)L_{-2}v,
 \qquad \|L_{-4}v\|^2=5c.
\]
Taking their vacuum pairings gives the matrix.
Its leading principal minors are positive when $c>0$.
\end{proof}

Assume $c>0$, and let $X_2,X_4$ be independent centered circular complex Gaussian
variables with variances $c/2$ and $5c$.
In their holomorphic polynomial space, the basis $X_4,X_2^2$
has Gram matrix $\operatorname{diag}(5c,c^2/2)$.
The finite-degree isometry fixing the first vector is
\begin{equation}\label{mht:eq:gravity-gram-isometry}
 L_{-4}v\longmapsto X_4,
 \qquad
 L_{-2}^2v\longmapsto
     \frac35X_4+\sqrt{1+\frac{22}{5c}}\,X_2^2.
\end{equation}

The coefficient $3/5$ gives the required cross term $3c$.
The remaining squared norm is
$c^2/2+4c-9c/5=(c^2/2)(1+22/(5c))$.
The assignment sending each Virasoro monomial directly to the
corresponding Gaussian monomial misses this cross term.
Both displayed forms are positive, and
\eqref{mht:eq:gravity-gram-isometry} corrects that finite assignment.
It does not yet intertwine the full Virasoro actions.

A physical unitary preserving the vacuum and conjugation must
preserve the Gram forms of its marked stress descendants.
Such a comparison uses the actual gauge quotient and real
structure of each state space. A holomorphic bilinear covariance
alone does not specify its Hermitian norm.

Direct Wick rotation also has a domain.
For instance,
\[
 u(t,\theta)=\sum_{n\ge1}e^{-\sqrt n}e^{-int+in\theta}
\]
is smooth for real $t,\theta$, since every polynomially weighted
coefficient sequence is summable.
At $t=i\tau$ with $\tau>0$, its Fourier coefficients become
$e^{n\tau-\sqrt n}$.
They exceed every polynomial bound, so they are not the
coefficients of a distribution on the circle.
For the opposite sign of imaginary time, the opposite frequency
sector gives the same obstruction.
Thus direct evaluation cannot rotate arbitrary smooth solutions.

By contrast, coefficients bounded by $Ce^{-\sigma n}$ give a
holomorphic continuation for $|\operatorname{Im}t|<\sigma$.
Finite Fourier sums are included in this analytic domain.
A comparison outside it requires another construction, such as
a boundary Green operator with specified closed domains.
The domain obstruction above does not decide whether such a
physical comparison exists.

\section{The Cardy expression and the BTZ area}

Suppose an independently defined boundary theory has the Cardy
asymptotic on a specified high-energy domain.
Write its cylinder energies as
$H_L=L_0-c_L/24$ and $H_R=\bar L_0-c_R/24$.
The leading entropy expression is
\[
 S_C=2\pi\sqrt{c_LH_L/6}+2\pi\sqrt{c_RH_R/6}.
\]
In a modular derivation, vacuum dominance after the modular
transformation and control of the inverse Laplace asymptotic
are part of this assumption.

For a nonextremal BTZ geometry, choose $G,\ell>0$ and
$r_+>|r_-|$. In the gravitational normalization, let
\begin{equation}\label{mht:eq:gravity-btz-charges}
 \begin{gathered}
 c_L=c_R=c=\frac{3\ell}{2G},\qquad
 M=\frac{r_+^2+r_-^2}{8G\ell^2},\qquad
 J=\frac{r_+r_-}{4G\ell},\\
 H_L=\frac{\ell M+J}{2}
     =\frac{(r_++r_-)^2}{16G\ell},\qquad
 H_R=\frac{\ell M-J}{2}
     =\frac{(r_+-r_-)^2}{16G\ell}.
 \end{gathered}
\end{equation}
These are cylinder charges. The vacuum shift must be retained
when comparing them with plane conformal weights.

\begin{proposition}\label{mht:prop:gravity-cardy-area}
Substitution of \eqref{mht:eq:gravity-btz-charges} into $S_C$ gives
\[
 S_C=\frac{2\pi r_+}{4G}.
\]
\end{proposition}

\begin{proof}
The two terms are
$\pi(r_++r_-)/(4G)$ and $\pi(r_+-r_-)/(4G)$.
Their sum gives the formula. The inequalities on the radii fix
the signs of the square roots.
\end{proof}

The right-hand side is the horizon circumference divided by $4G$.
The equality identifies the leading Cardy expression with the
area term wherever the assumed asymptotic applies.
It does not determine the finite-energy degeneracies or establish
the existence of the boundary theory.\footnote{Strominger,
\emph{Black hole entropy from near-horizon microstates},
Sections~3 and~5, keeps the cylinder shift and states the
high-energy regime and the existence assumption,
\url{https://arxiv.org/abs/hep-th/9712251}.}

\section{Marked classical lines}

The projective matrix connection also gives an explicit classical
spin-one-half line. Let $q$ be smooth on a compact interval, set
$a(t)=\left(\begin{smallmatrix}0&q(t)\\1&0\end{smallmatrix}\right)$,
and let $U(y,x)$ solve
\begin{equation}\label{mht:eq:gravity-marked-transport}
 \partial_yU(y,x)=-a(y)U(y,x),\qquad U(x,x)=I.
\end{equation}
For a column solution, its second component satisfies
$\psi_2''=q\psi_2$.
The endpoint bases are retained, so $U$ is a matrix rather than
its trace.

\begin{proposition}\label{mht:prop:gravity-marked-lines}
The transport has determinant one and satisfies
\begin{equation}\label{mht:eq:gravity-line-sewing}
 U(z,x)=U(z,y)U(y,x).
\end{equation}
For an arbitrary coefficient variation,
\begin{equation}\label{mht:eq:gravity-line-variation}
 \delta U(y,x)
 =-\int_x^yU(y,t)\,\delta a(t)\,U(t,x)\,dt.
\end{equation}
For $\delta a=\lambda'+[a,\lambda]$, this becomes
\begin{equation}\label{mht:eq:gravity-line-endpoints}
 \delta U(y,x)=-\lambda(y)U(y,x)+U(y,x)\lambda(x).
\end{equation}
\end{proposition}

\begin{proof}
Iterating the integral equation for
\eqref{mht:eq:gravity-marked-transport} gives the ordered series
\[
 I+\sum_{n\ge1}(-1)^n
 \int_{x<t_1<\cdots<t_n<y}
       a(t_n)\cdots a(t_1)\,dt_1\cdots dt_n.
\]
Its $n$th norm is at most
$(|y-x|\sup\|a\|)^n/n!$, so it converges uniformly.
Uniqueness of the initial-value problem gives
\eqref{mht:eq:gravity-line-sewing}.
Jacobi's determinant identity gives
$\partial_y\det U=-\operatorname{tr}(a)\det U=0$.

Differentiate the integral equation with respect to the
coefficient variation to obtain \eqref{mht:eq:gravity-line-variation}.
For a gauge variation the integrand is
\[
 \partial_t\bigl(U(y,t)\lambda(t)U(t,x)\bigr).
\]
Its integral gives \eqref{mht:eq:gravity-line-endpoints}.
\end{proof}

The same identities give a classical ghost differential.
For an odd matrix ghost $c$, set
\[
 Qa=c'+[a,c],\qquad Qc=-c^2,
 \qquad QU(y,x)=-c(y)U(y,x)+U(y,x)c(x).
\]
The graded Jacobi identity gives $Q^2a=Q^2c=0$.
Using the Leibniz sign for $Q(cU)$ gives $Q^2U=0$.

Applying $Q$ to \eqref{mht:eq:gravity-line-sewing} cancels the
two ghost actions at the common endpoint.
Thus marked parallel transport is compatible with the
classical gauge complex and ordered sewing.
The identity matrix is its interval-composition unit.
Comparison with a local quantum boundary vacuum requires the
corresponding unit action map.

\section{The physical comparison maps}

The classical constructions above leave distinct quantum questions.
Their source and target complexes must be defined with their
own actions, support conditions, completions and states.
The following problem specifies the maps needed to join them.

\begin{problem}\label{mht:prob:gravity-native-maps}
Construct the following comparisons on matched physical domains.
\begin{enumerate}
\item For the sourced solid-torus Gaussian, construct its
renormalized determinant or Pfaffian line, connection and phase.
Compare them with an independently normalized boundary preparation,
preserving source derivatives, disk-cut Green data and cap sewing.
The permitted finite bulk, boundary and corner corrections must
preserve the declared source identity.

\item Construct an insertion from the affine-current and ghost
complex into the constrained native observable complex.
Require the chain identity
$Q_{\mathrm{phys}}\mathcal I=\mathcal I Q_{\mathrm{DS}}$,
including composite current products, their simultaneous
cutoff/collision limits, translation and the vacuum action.
Derive the relation between the affine level and gravitational
coupling from this map and the physical normalization.

\item Construct the physical Euclidean cut-state quotient,
its Lorentzian comparison and its conjugation.
On this quotient, construct both Virasoro actions and the
marked stress descendants with their Hermitian forms.
Compare the normalized cylinders and the sewing pairing.

\item Construct the coincident-junction products of physical
cycle operators on a common nonzero state domain.
Retain their determinant-derived framing, monodromy and local
vacuum action, and determine the unreduced Weyl parameter there.
Formal power-series constructions and fixed numerical levels
require their own comparison of domains.

\item Define independently an effective spin-one-half bulk-line
complex $\mathscr L(y,x)$ and a marked boundary-operation complex
$\operatorname{Op}^{\mathrm{marked}}_\partial(y,x)$.
Construct a degree-zero comparison
\[
 \beta_{y,x}:\mathscr L(y,x)\longrightarrow
 \operatorname{Op}^{\mathrm{marked}}_\partial(y,x),
\]
satisfying the chain identity and both endpoint identities:
\[
 \begin{aligned}
 Q_\partial\beta&=\beta Q_{\mathrm{eff}},\\
 \beta(a_{\mathrm{out}}u)
       &=\rho_{\mathrm{out}}(a_{\mathrm{out}})\beta(u),\\
 \beta(u a_{\mathrm{in}})
       &=\beta(u)\rho_{\mathrm{in}}(a_{\mathrm{in}}).
 \end{aligned}
\]
Here $\rho_{\mathrm{out}},\rho_{\mathrm{in}}$ are the independently
specified endpoint actions.
Require these identities and ordered sewing at every formal order.
The selected renormalized staple $\mathcal W_t^{\mathrm{ren}}$
must be a cocycle at every
positive depth, with its independently normalized quantum Hill
field as boundary image.
Construct a degree-minus-one depth homotopy $\mathcal H_t$ with
$\partial_t\mathcal W_t^{\mathrm{ren}}
=-Q_{\mathrm{eff}}\mathcal H_t$ and a compatible marked image.

The construction must retain normal heat profiles, ghosts,
antifields and common subgraph restrictions.
In particular, construct the quadratic coupling coefficient
with both endpoint actions and sewing.
\end{enumerate}
\end{problem}

An obstruction to one of these maps must survive the complete
admissible counterterm complex on its specified carrier.
The matrix, Gaussian and finite-state examples distinguish
particular constructions and their domains.
They supply exact conditions that a physical comparison must
preserve, while its existence remains a separate mathematical
question.
